\section{Online Learning Problem}
\label{sec:online_learning}

In the online learning setting, the receiver does not know the joint distribution of $(U,X,Y)$, while the channel noise parameter $q \in [0,1]$ is assumed known. Instead of observing a static offline training dataset, the receiver sequentially interacts with the unknown environment over rounds $t = 1, 2, \ldots, T$. At each round $t$, a fresh sender sample $(X_t, Y_t, U_t)$ is first drawn from the underlying environment distribution $P_{(U,X,Y)}$. The receiver then selects a disclosure regime $m_t \in \mathcal{M} := \{V, S\}$ and a regression parameter $b_t \in \mathcal{R} := \{b \in \mathbb{R}^d : \|b\| \le 1\}$ (the compact strategy space defined in Section~\ref{sec:offline_learning}), forming the joint action $a_t = (m_t, b_t) \in \mathcal{A} := \mathcal{M} \times \mathcal{R}$. At the end of round $t$, the receiver observes scalar bandit loss feedback $Z_t = \ell(m_t, b_t, X_t, Y_t, U_t)$, which is the least-squares loss for the realized $(U_t, X_t, Y_t)$ under model $b_t$ and disclosure regime $m_t$. The receiver then updates both its chosen regime and parameter vector to minimize cumulative loss, aiming to match the performance of the optimal fixed regime and regression model in hindsight.

We pose this as a continuous bandit feedback problem and approach it through Gaussian Process Thompson Sampling (GP-TS) algorithm as proposed in Algorithm 2 in \citet{chowdhury2017kernelized}. The receiver aims to optimize the unknown expected risk function $f(a)$ over the joint action space $\mathcal{A} = \mathcal{M} \times \mathcal{R}$, defined as
\begin{equation}
\label{eq:expected_risk_definition}
f(a) := \mathbb{E}_{(X,Y,U)}[\ell(m,b,X,Y,U)] = \begin{cases} \RV(b) & \text{if } m = \text{Vanilla}, \\ \RS(b) & \text{if } m = \text{Strategic}. \end{cases}
\end{equation}
To address this, the receiver places a Gaussian Process prior over candidate loss functions, initialized with prior mean $\mu_0(a) = 0$ and prior covariance kernel $k_0(a, a') = k(a, a')$. For joint actions $a = (m, b)$ and $a' = (m', b')$, the kernel $k_0(a, a') = k(a, a')$ takes the Matérn-$5/2$ functional form:
\begin{equation}
\label{eq:matern_kernel_def}
k(a, a') := \mathbb{I}\{m = m'\}\left(1 + \frac{\sqrt{5}\|b - b'\|}{l} + \frac{5\|b - b'\|^2}{3 l^2}\right) \exp\left( -\frac{\sqrt{5}\|b - b'\|}{l} \right),
\end{equation}
where $l > 0$ is the kernel lengthscale parameter. For the direct-sum kernel $k((m,b),(m',b')) = \mathbb{I}\{m=m'\}k_0(b,b')$, the joint RKHS norm satisfies $B_k \le \sqrt{2}\max\{B_{k,V}, B_{k,S}\}$. At each round $t$, the receiver maintains the updated posterior distribution $\operatorname{GP}(\mu_{t-1}, v_t^2 k_{t-1})$, where $v_t$ is an inflation factor given by $v_t = B_k + R\sqrt{2(\gamma_{t-1}+1+\log(2/\delta))}$ \citep{chowdhury2017kernelized}. The receiver draws a virtual objective function $\widetilde{f}_t \sim \operatorname{GP}(\mu_{t-1}, v_t^2 k_{t-1})$ and selects action $a_t = (m_t, b_t) \in \arg\min_{a \in \mathcal{A}_t} \widetilde{f}_t(a)$ over the decision set $\mathcal{A}_t = \{V, S\} \times \mathcal{R}_t$, where $\mathcal{R}_t$ is an $\epsilon_t$-discretization net of $\mathcal{R}$ of mesh size $\epsilon_t = \mathcal{O}(1/t^2)$ following \citet{chowdhury2017kernelized}.

Upon observing the scalar loss feedback $Z_t$, the receiver updates the GP posterior mean $\mu_t(a)$ and covariance kernel $k_t(a, a')$ according to the standard Gaussian process regression update rules \citep{chowdhury2017kernelized}:
\begin{align}
    \mu_t(a) &= \mathbf{k}_t(a)^\top (K_t + \lambda I_t)^{-1} \mathbf{Z}_t, \label{eq:gp_mean_update} \\
    k_t(a, a') &= k(a, a') - \mathbf{k}_t(a)^\top (K_t + \lambda I_t)^{-1} \mathbf{k}_t(a'), \label{eq:gp_kernel_update}
\end{align}
where $\mathbf{Z}_t = [Z_1, \ldots, Z_t]^\top \in \mathbb{R}^t$ is the vector of observed losses, $\mathbf{k}_t(a) = [k(a_1, a), \ldots, k(a_t, a)]^\top \in \mathbb{R}^t$ is the kernel cross-covariance vector between action $a$ and past actions $a_1, \ldots, a_t$, and $K_t \in \mathbb{R}^{t \times t}$ is the kernel Gram matrix with entries $[K_t]_{i,j} = k(a_i, a_j)$ for $1 \le i, j \le t$, with $\lambda > 0$ serving as the regularization parameter. This procedure is detailed in Algorithm~\ref{alg:kernel_ts}.

\begin{algorithm}[htbp]\index{Gaussian Process Thompson Sampling (GP-TS)|textbf}\index{Online learning problem|textbf}
    \caption{Gaussian Process Thompson Sampling (GP-TS)}
    \label{alg:kernel_ts}
    \small
    \begin{algorithmic}[1]
        \State \textbf{Input:} Matérn-$5/2$ covariance kernel $k_0(a,a') = k(a,a') = \mathbb{I}\{m = m'\} \left(1 + \frac{\sqrt{5}\|b - b'\|}{l} + \frac{5\|b - b'\|^2}{3 l^2}\right) \exp\left( -\frac{\sqrt{5}\|b - b'\|}{l} \right)$, regularization $\lambda > 0$, confidence $\delta \in (0,1)$.
        \State \textbf{Initialize:} Posterior mean $\mu_0(a) = 0$ and covariance $k_0(a,a') = k(a,a')$ for all $a \in \mathcal{A}$.
        \For{$t = 1, 2, \ldots, T$}
            \State Set inflation factor $v_t = B_k + R\sqrt{2(\gamma_{t-1}+1+\log(2/\delta))}$.
            \State Sample virtual objective function $\widetilde{f}_t \sim \operatorname{GP}(\mu_{t-1}, v_t^2 k_{t-1})$.
            \State Select action $a_t = (m_t, b_t) \in \arg\min_{a \in \mathcal{A}_t} \widetilde{f}_t(a)$ over decision set $\mathcal{A}_t = \{V, S\} \times \mathcal{R}_t$.
            \State Observe loss feedback $Z_t = \ell(m_t, b_t, X_t, Y_t, U_t)$.
            \State Update posterior mean $\mu_t(a)$ and covariance $k_t(a,a')$ via \eqref{eq:gp_mean_update} and \eqref{eq:gp_kernel_update}.
        \EndFor
    \end{algorithmic}
\end{algorithm}

We measure the performance of Algorithm~\ref{alg:kernel_ts} using cumulative regret relative to the minimal expected loss achievable in hindsight, namely $\min\{\RS^\star, \RV^\star\}$:
\begin{equation}
\label{eq:online_regret_definition}\index{Cumulative regret|textbf}
    \operatorname{Reg}_T := \sum_{t=1}^T \bigl(f(a_t) - f(a^\star)\bigr), \qquad f(a^\star) = \min\{\RS^\star, \RV^\star\},
\end{equation}
where $a^\star \in \arg\min_{a \in \mathcal{A}} f(a)$.

To establish a regret guarantee for GP-TS in our setting, we define the following class of regular environments.

\begin{definition}[Sobolev Environments]\label{def:sobolev_risk}\index{Sobolev environments|textbf}\index{Sobolev space|textbf}
Let $\mathcal{P}_{\text{Sob}}$ denote the set of environment joint distributions $(U,X,Y)$ such that the expected risk functions $L(m,b)$ belong to the Sobolev space $H^{5/2 + d/2}(\Rspace)$.
\end{definition}

We next show that $\mathcal{P}_{\text{Sob}}$ contains all environments with a smooth joint density $p_{U,X,Y}(u,x,y)$ on compact support, and is therefore non-empty.

\begin{lemma}[Smooth Environments]
\label{lem:sobolev_membership}
Suppose the environment variables $(U, X, Y)$ have a smooth, bounded joint density $p_{U,X,Y}(u,x,y)$ with compact support, and the receiver's action space $\mathcal{R} \subset \{b \in \mathbb{R}^d : \|b\|_2 \ge b_{\min}\}$ is a compact domain bounded away from the origin for some $b_{\min} > 0$. Then $L(m, b) \in H^{5/2 + d/2}(\mathcal{R})$, and the environment is contained in $\mathcal{P}_{\text{Sob}}$.
\end{lemma}
\begin{proof} The proof is given in Section~\ref{sec:proofs_ch4} (Proof~\ref{proof:sobolev_membership}).
\end{proof}

We now state the theorem which gives the regret of Algorithm~\ref{alg:kernel_ts}.

\begin{theorem}[Regret Bound for Strategic GP-TS]\label{thm:kernel_online_regret}\index{Mat\'ern kernel|textbf}\index{Cumulative regret bound|textbf}
For any environment distribution in $\mathcal{P}_{\text{Sob}}$, let the receiver run GP-TS (Algorithm~\ref{alg:kernel_ts}) with a Matérn-$\nu$ kernel (for smoothness parameter $\nu > 0$) over the joint action space $\mathcal{A} = \mathcal{M} \times \mathcal{R} \subset \{V, S\} \times \mathbb{R}^d$. Then, with probability at least $1 - \delta$, the cumulative regret satisfies
\begin{equation}
\label{eq:kernel_online_regret_bound}
    \operatorname{Reg}_T = \widetilde{\mathcal{O}}\left( T^{\frac{2\nu+3d}{4\nu+2d}} \right).
\end{equation}
Specifically, substituting the practical Matérn-$5/2$ kernel ($\nu=5/2$), the bound yields $\widetilde{\mathcal{O}}\bigl( T^{\frac{5+3d}{10+2d}} \bigr)$, which is sublinear ($\operatorname{Reg}_T/T \to 0$) for dimensions $d \in \{1, 2, 3, 4\}$.
\end{theorem}
\begin{proof}[Proof Outline]
The proof of Theorem~\ref{thm:kernel_online_regret} applies the kernelized Thompson Sampling regret bound of \citet[Theorem 4]{chowdhury2017kernelized}, bounding the maximum information gain by $\gamma_T = \mathcal{O}\left(T^{\frac{d}{5+d}} (\log T)^{\frac{5}{5+d}}\right)$ \citep[Table 1, Theorem 1]{vakili2021information}. The complete proof is given in Section~\ref{sec:proofs_ch4} (Proof~\ref{proof:kernel_online_regret}).
\end{proof}

\paragraph{Simulation Results for Online Learning.}

To evaluate Algorithm~\ref{alg:kernel_ts} and empirically validate the cumulative regret bounds from Theorem~\ref{thm:kernel_online_regret}, we simulate the online learning problem across feature dimensions $d \in \{1, 2\}$.\coderef{subsec:code_online_learning} For each evaluation environment, the data generation process draws joint feature vectors, outcomes, and sender preferences $(X, Y, U) \in \mathbb{R}^{d+2}$ from a 2-component Gaussian Mixture Model (GMM) with mixture weights $\pi_1 = 0.6, \pi_2 = 0.4$. Component 1 represents a sub-population with mean vector $\mu_1 = (+0.8 \mathbf{1}/\sqrt{d}, 1.0, 1.0)^\top$, while Component 2 represents a sub-population with mean vector $\mu_2 = (-0.8 \mathbf{1}/\sqrt{d}, -0.8, -1.0)^\top$. Both components share a symmetric positive-definite covariance matrix $\Sigma \in \mathbb{R}^{(d+2) \times (d+2)}$ specified by feature covariance $\Sigma_X = 0.64 I_d$, outcome variance $\mathrm{Var}(Y) = 0.6$, preference variance $\mathrm{Var}(U) = 0.5$, and cross-covariances $\mathrm{Cov}(X,Y) = \frac{0.5}{\sqrt{d}}\mathbf{1}$, $\mathrm{Cov}(X,U) = \frac{0.3}{\sqrt{d}}\mathbf{1}$, and $\mathrm{Cov}(Y,U) = 0.3$. The channel noise parameter is fixed at $q = 0.5$.

In each independent trial, the receiver sequentially interacts with the environment over $T = 300$ rounds using Algorithm~\ref{alg:kernel_ts} configured with a Matérn kernel (approximated practically using a Squared-Exponential kernel with lengthscale $l = 0.7$) and regularization $\lambda = 0.1$. At each round $t$, a sample function $\widetilde{f}_t$ is drawn from the updated GP posterior to select action $a_t = (m_t, b_t) \in \mathcal{A}$. We average the cumulative regret $\operatorname{Reg}_T$ across $M = 10$ independent Monte Carlo simulation runs.

The empirical findings are presented in Figure~\ref{fig:online_regret_validation}. As predicted by Theorem~\ref{thm:kernel_online_regret}, the empirical cumulative regret exhibits a sublinear growth rate matching the theoretical $\widetilde{\mathcal{O}}\left(T^{\frac{5+3d}{10+2d}}\right)$ scaling (which instantiates the general Matérn regret bound $\widetilde{\mathcal{O}}\bigl(T^{\frac{2\nu+3d}{4\nu+2d}}\bigr)$ for smoothness parameter $\nu = 5/2$), with shaded bands showing $\pm 1$ standard error across $M=10$ independent Monte Carlo trials. Furthermore, comparing $d=1$ and $d=2$ confirms that higher feature dimensionality leads to a larger regret growth rate due to the dimensional dependence, corroborating our theoretical analysis for environments with smooth joint densities $p_{U,X,Y}$ on compact support.

\begin{figure}[htbp]
    \centering
    \includegraphics[width=0.72\textwidth]{online_learning_regret.pdf}
    \caption{Empirical cumulative regret $\operatorname{Reg}_T$ versus time horizon $T$ for $d=1$ and $d=2$ (solid lines depict mean cumulative regret, shaded bands indicate $\pm 1$ standard error across $M=10$ trials).}
    \label{fig:online_regret_validation}
\end{figure}

% Applied Fix: Restricted smooth density lemma to domains strictly separated from b=0 and specified nu=5/2.
